\documentclass[../../main.tex]{subfiles}
\begin{document}

{\bf [解]}
新しく
\begin{align}
  f(n) 
   &=\dfrac{{}_{3n}C_n}{{}_{2n}C_n}=\dfrac{(3n)!}{n!(2n)!}\cdot\dfrac{n!n!}{(2n)!}=\dfrac{(3n)!\,n!}{(2n)!(2n)!}
\end{align}
とおくと$\left(f(n)\right)^{\frac{1}{n}}$の$n\to\infty$の極限を求めれば良い．
$f(n)>0$より$g(n)=\log \left(f(n)^{\frac{1}{n}}\right)$とおくと
\begin{align}
  g(n)
   &= \dfrac{1}{n}\log \left(\dfrac{(3n)!\,n!}{(2n)!(2n)!}\right) \\
   &= \dfrac{1}{n}\left(\log (3n)! + \log n! - 2\log(2n)!\right) \\
   &= \dfrac{1}{n}\left(\sum_{k=1}^{3n}\log k + \sum_{k=1}^{n}\log k  - 2\sum_{k=1}^{2n}\log k \right) \\
   &= \dfrac{1}{n}\left(\sum_{k=2n+1}^{3n}\log k - \sum_{k=n+1}^{2n}\log k \right) \\
\end{align}
ここで$\log n - \log n = 0 $を間に挟むと
\begin{align}
 g(n) 
   &= \dfrac{1}{n}\left(\sum_{k=2n+1}^{3n}\log \dfrac{k}{n} - \sum_{k=n+1}^{2n}\log \dfrac{k}{n} \right) \label{eq:1}
\end{align}
と書ける．ここで区分求積法より
\begin{align}
 \lim_{n\to\infty}\dfrac{1}{n}\sum_{k=n+1}^{2n}\log\dfrac{k}{n} = \int_1^2\log x\,dx&=2(\log2+1)-1 \\
 \lim_{n\to\infty}\dfrac{1}{n}\sum_{k=2n+1}^{3n}\log\dfrac{k}{n} = \int_2^3\log x\,dx&=3(\log3+1)-2(\log2+1)
\end{align}
だから\cref{eq:1}に代入して
\begin{align}
 &\lim_{n\to\infty}g(n) = 3\log3+4-4\log2-4=3\log3-4\log2=\log\dfrac{27}{16} \\
\therefore 
 &\lim_{n\to\infty} \log \left(f(n)^{\frac{1}{n}}\right) = \log\dfrac{27}{16}
\end{align}
である．$\log x$の連続性から
\begin{align}
 \lim_{n\to\infty} \left(f(n)^{\frac{1}{n}}\right) = \dfrac{27}{16}
\end{align}
が求める極限値である．

\newpage
\end{document}
